\begin{examplebox}
\textbf{\textcolor{CExample}{例 1（基础）}}

\textbf{\textcolor{CExample}{题目：}} 已知三棱锥 $S$-$ABC$ 中，$SA=SB=SC=5$，底面 $\triangle ABC$ 是边长为 $6$ 的正三角形。求该三棱锥的外接球半径 $R$。

\textbf{\textcolor{CExample}{解题步骤：}}
\begin{description}[style=nextline, leftmargin=2.8em, labelsep=0.5em, itemsep=0.45em]
    \item[\textbf{第一步（侧棱判断）：}] 由 $SA=SB=SC$ 得 $S$ 在底面的射影为 $\triangle ABC$ 的外心 $O$，即正三角形的中心，且 $SO\perp$ 底面。
    \par\textbf{理由：} 在非退化三棱锥中，侧棱相等 $\Rightarrow$ 顶点射影为底面外心。

    \item[\textbf{第二步（求量计算）：}] 正三角形边长为 $6$，外接圆半径 $r = \dfrac{6}{\sqrt{3}} = 2\sqrt{3}$；在 $\mathrm{Rt}\triangle SOA$ 中，$h = \sqrt{SA^{2}-r^{2}} = \sqrt{25-12}=\sqrt{13}$。
    \[
    r = \frac{6}{\sqrt{3}} = 2\sqrt{3},\qquad h = \sqrt{5^{2}-(2\sqrt{3})^{2}} = \sqrt{13}.
    \]
    \par\textbf{理由：} 正三角形外接圆公式及勾股定理。

    \item[\textbf{第三步（应用结论）：}] 利用变形公式 $R = \dfrac{h^{2}+r^{2}}{2h}$ 得
    \[
    R = \frac{13+12}{2\sqrt{13}} = \frac{25}{2\sqrt{13}} = \frac{25\sqrt{13}}{26}.
    \]
    \par\textbf{理由：} 由 $R^{2}=r^{2}+OQ^{2}$ 且 $R=|h-OQ|$ 推导得到。
\end{description}

\textbf{\textcolor{CExample}{关键结论：}} \textbf{$R = \dfrac{25\sqrt{13}}{26}$}.

\medskip
\textbf{\textcolor{CExample}{例 2（稍变形）}}

\textbf{\textcolor{CExample}{题目：}} 三棱锥 $S$-$ABC$ 中，$SA=SB=SC=\sqrt{41}$，底面 $\triangle ABC$ 中 $\angle A=90^{\circ}$，$AB=3$，$AC=4$。求外接球半径 $R$。

\textbf{\textcolor{CExample}{解题步骤：}}
\begin{description}[style=nextline, leftmargin=2.8em, labelsep=0.5em, itemsep=0.45em]
    \item[\textbf{第一步（侧棱判断）：}] 由 $SA=SB=SC$ 得 $S$ 在底面射影为 $\triangle ABC$ 的外心 $O$。直角三角形外心为斜边中点，故 $O$ 为 $BC$ 中点，$SO\perp$ 底面。
    \par\textbf{理由：} 在非退化三棱锥中，侧棱相等等价于顶点射影为底面外心。

    \item[\textbf{第二步（计算 $r$ 和 $h$）：}] \ $\displaystyle BC = \sqrt{AB^{2}+AC^{2}} = 5$，所以 $r = \dfrac{BC}{2} = \dfrac{5}{2}$。在 $\mathrm{Rt}\triangle SOA$ 中，
    \[
    h = \sqrt{SA^{2}-r^{2}} = \sqrt{41-\frac{25}{4}} = \sqrt{\frac{139}{4}} = \frac{\sqrt{139}}{2}.
    \]
    \par\textbf{理由：} 勾股定理与直角三角形外接圆性质。

    \item[\textbf{第三步（应用结论）：}] 
    \[
    R = \frac{h^{2}+r^{2}}{2h} = \frac{\frac{139}{4}+\frac{25}{4}}{2\cdot\frac{\sqrt{139}}{2}} = \frac{164}{4\sqrt{139}} = \frac{41}{\sqrt{139}} = \frac{41\sqrt{139}}{139}.
    \]
    \par\textbf{理由：} 半径公式。
\end{description}

\textbf{\textcolor{CExample}{关键结论：}} \textbf{$R = \dfrac{41\sqrt{139}}{139}$}.

\medskip
\textbf{\textcolor{CExample}{例 3（综合应用）}}

\textbf{\textcolor{CExample}{题目：}} 三棱锥 $S$-$ABC$ 中，$SA=SB=SC$，底面 $\triangle ABC$ 满足 $AB=AC=4$，$BC=6$，顶点 $S$ 到底面 $ABC$ 的距离 $h=2$。求外接球半径 $R$。

\textbf{\textcolor{CExample}{解题步骤：}}
\begin{description}[style=nextline, leftmargin=2.8em, labelsep=0.5em, itemsep=0.45em]
    \item[\textbf{第一步（侧棱判断）：}] 由 $SA=SB=SC$ 得 $S$ 在底面射影为 $\triangle ABC$ 的外心 $O$，$SO\perp$ 底面。由 $\triangle ABC$ 等腰，$O$ 在底边的中垂线上。
    \par\textbf{理由：} 在非退化三棱锥中，侧棱相等等价于顶点射影为底面外心。

    \item[\textbf{第二步（求底面外接圆半径 $r$）：}] 等腰 $\triangle ABC$ 底边 $BC=6$，腰 $AB=AC=4$，底边上的高 $h_{\triangle}= \sqrt{4^{2}-3^{2}} = \sqrt{7}$，面积 $S_{\triangle}= \frac{1}{2}\cdot6\cdot\sqrt{7}=3\sqrt{7}$。由正弦定理推论 $r = \dfrac{AB\cdot AC\cdot BC}{4S_{\triangle}}$，得
    \[
    r = \frac{4\cdot4\cdot6}{4\cdot3\sqrt{7}} = \frac{96}{12\sqrt{7}} = \frac{8}{\sqrt{7}} = \frac{8\sqrt{7}}{7}.
    \]
    \par\textbf{理由：} 三角形外接圆半径公式 $r = \dfrac{abc}{4S}$。

    \item[\textbf{第三步（应用结论）：}] 直接代入半径公式 $R = \dfrac{h^{2}+r^{2}}{2h}$：
    \[
    R = \frac{2^{2} + \left(\frac{8\sqrt{7}}{7}\right)^{2}}{2\cdot2} = \frac{4 + \frac{64}{7}}{4} = \frac{\frac{28+64}{7}}{4} = \frac{92}{28} = \frac{23}{7}.
    \]
    \par\textbf{理由：} 外接球半径公式。
\end{description}

\textbf{\textcolor{CExample}{关键结论：}} \textbf{$R = \dfrac{23}{7}$}.

\end{examplebox}